\documentclass[
  article,
  12pt,
  oneside,
  english,
  brazil,
	%sumario=tradicional
  sumario=abnt-6027-2012,
]{abntex2}

\usepackage[alf]{abntex2cite}
\renewcommand{\bibsection}{
  %\section*{Referências bibliográficas}
  \section{Referências bibliográficas}
  %\phantomsection
  %\addcontentsline{toc}{section}{Referências bibliográficas}
}

% Set all sections to break pages if they cant fit properly
\let\oldsection\section
\renewcommand{\section}{\filbreak\oldsection}
% And same for chapters
\let\oldchapter\chapter
\renewcommand{\chapter}{\filbreak\oldchapter}

\usepackage{indentfirst}	
\usepackage{graphicx}
\usepackage{wrapfig}
\usepackage{microtype}
\usepackage[utf8]{inputenc}
\usepackage{xcolor}
\usepackage{subcaption}

% Tikz config
\usepackage{tikz}
\usepackage{tikz-3dplot}
\usetikzlibrary{patterns, patterns.meta}

%BEGIN Coding style
\usepackage{listings}
\definecolor{codegreen}{RGB}{64, 160, 43}
\definecolor{subtext0}{RGB}{108, 111, 133}
\definecolor{codemauve}{RGB}{136, 57, 239}
\definecolor{codepeach}{RGB}{254, 100, 11}
\definecolor{codebg}{RGB}{255, 255, 255}
\lstdefinestyle{mystyle}{
  backgroundcolor=\color{codebg},
  commentstyle=\color{subtext0},
  keywordstyle=\color{codemauve},
  numberstyle=\tiny\color{subtext0},
  stringstyle=\color{codegreen},
  basicstyle=\ttfamily\footnotesize\linespread{0.5}\selectfont,
  breakatwhitespace=false,
  breaklines=true,
  captionpos=b,
  keepspaces=true,
  numbers=left,
  numbersep=3pt,
  showspaces=false,
  showstringspaces=false,
  showtabs=false,
  tabsize=4,
}
\lstset{style=mystyle}
%END Coding style

\newcounter{generatedcode}
\setcounter{generatedcode}{0}
\newcommand{\getcode}[2]{
  \immediate\addtocounter{generatedcode}{1}
  \immediate\write18{./generate_code #1 #2 > /tmp/tempcode-\thegeneratedcode.txt}
  \immediate\lstinputlisting[language=Python]{/tmp/tempcode-\thegeneratedcode.txt}
  \immediate\write18{rm /tmp/tempcode-\thegeneratedcode.txt}
}


\usetikzlibrary{shapes, arrows, positioning}
\tikzstyle{terminator} = [rectangle, draw, text centered, rounded corners, minimum height=2em]
\tikzstyle{process} = [rectangle, draw, text centered, minimum height=2em]
\tikzstyle{decision} = [diamond, draw, text centered, minimum height=2em, minimum width=3em, aspect=2]
\tikzstyle{data} = [trapezium, draw, text centered, trapezium left angle=60, trapezium right angle=120, minimum height=2em]
\tikzstyle{connector} = [draw, -latex']

%\newcommand{\todo}[1]{\textcolor{red}{WIP #1}}
\newenvironment{wip}{
  \color{violet}
}{
  \ignorespacesafterend
}
\newenvironment{todo}{
  \color{red}
}{
  \ignorespacesafterend
}

\titulo{Manual: P.H.A.A.S.T}
\autor{Daniel Haas}
\data{2025}
\local{Universidade Federal do Rio de Janeiro}
\orientador{Ricardo Rodrigues de Oliveira Junior}
\tipotrabalho{Tutorial}

\begin{document}

\imprimircapa

\section{Usage}

PHAAST has two main non-optional and positional arguments, \lstinline{STOICHIOMETRY}, which defines the atoms that will compose the searched minima, and \lstinline{OUTPUT_PREFIX}, a starting name given to the resulting minima that will be stored in xyz format.
The command can be summarized in:

\begin{lstlisting}
  python -m phaast [OPTIONS] STOICHIOMETRY OUTPUT_PREFIX  
\end{lstlisting}

\subsection{Options}
All other option arguments can be found on the help section, which is shown with the use of the option \lstinline{-h}:

\begin{lstlisting}
-h, --help            show this help message and exit

--return_number RETURN_NUMBER
  Number of least energy molecules to return at the end of the algorithm (default = 50)

-pop, --population_size POPULATION_SIZE
  Number of molecules in the population (default = 5000)

-c, --charge CHARGE   Structure's charge (default = 0)

-t, --threads THREADS
  Number of threads to be used by the algorithm (default = number of threads on cpu)

-xtbt, --xtb-threads XTB_THREADS
  Number of threads to be used by each -t xtb instance (default = 4)

--xtb_path XTB_PATH
  Path to be used for the xtb binary (default = xtb)

-et, --energy_threshold ENERGY_THRESHOLD
  Maximum energy difference [in hartree] so molecules are considered alike (default = 0.001)

-gt, --geometry_threshold GEOMETRY_THRESHOLD
  Maximum geometry difference so molecules are considered the same so one of them is discarded [must be a value between 0 and 1] (default = 0.87)

--mut_disp_w MUT_DISP_W
  Mutation displacement weight (default = 3)

--mut_perm_w MUT_PERM_W
  Mutation permutation weight (default = 2)

--mut_twist_w MUT_TWIST_W
  Mutation twist weight (default = 1)

--crov_w CROV_W
  Crossing-over weight (default = 1)

--migr_w MIGR_W
  Migrator weight (default = 1)

--mut_disp_min MUT_DISP_MIN
  Minimum distance (A) to be used in mutation of atomic displacement (default = 1.3)

--mut_disp_max MUT_DISP_MAX
  Maximum distance (A) to be used in mutation of atomic displacement (default = 2.3)

--mut_disp_num MUT_DISP_NUM
  Number of atoms to move in mutations of atomic displacement (default = 1)

--mut_perm_num MUT_PERM_NUM
  Number of permutations in mutations of that type (default = Half the number of atoms)

--mut_twist_min_angle MUT_TWIST_MIN_ANGLE
  Minimum angle to be used in twist mutations (default = 90)

--mut_twist_max_angle MUT_TWIST_MAX_ANGLE
  Maximum angle to be used in twist mutations (default = 270)

-loops, --end_loop_number END_LOOP_NUMBER
  Number of loops the algorithm will run with the least energy molecule until a new one is found, terminating it (default = 11)

--comparison_algorithm COMPARISON_ALGORITHM
  Algorithm to be used for structures comparison, the recomendes usage is:
    - "grigoryan_springborg" for clusters;
    - "bonding_length" for organic and general shaped structures (check --comparison_bonding_tolerance too);
    (default = grigoryan_springborg)

--comparison_bonding_tolerance COMPARISON_BONDING_TOLERANCE
  Distance tolerable so two atoms can be considered bonded within the molecule comparison algorithm and bonding tests

--remove_unbonded REMOVE_UNBONDED
  Choose to not remove unbonded structures from genetic algorithm population
    - "always": Remove unbonded in every loop;
    - "final": Remove only from final population to not contaminate results;
    - "never": Do not remove;
    (default = always)

-v, --verbose
\end{lstlisting}

\subsubsection{Return number}

Changing the \lstinline{return_number} will return more or less structures at the end of the algorithm, there's no additional overhead to that, since the structures have already been optimized up to that point.
It is useful when the calculator used to optimize the structures is not so accurate on the energy values and a better optimization (on a more refined functional/basisset) is needed at the end of the algorithm, hence justifying a smaller return number.

\subsubsection{Population size}
Despite how it sounds, this argument doesn't define a cap or a direct control on the population's size, it actually defines how many structures will be generated in each generation, respecting the \textbf{Weights} defined at the optional arguments.

\subsubsection{Threads}
Number of processes (or threads if you are using a GIL-disabled python instance) used for paralelizing the expensive tasks on the genetic algorithm.

\subsubsection{Weights}
All the \lstinline{--*_w} arguments define the proportions of how each operation will be applied.
For instance, \lstinline{--mut_disp_w}, \lstinline{--mut_perm_w} and \lstinline{--mut_twist_w} define all mutation weights, if they're the same value, all mutations will occur in the same proportion.



\end{document}
